\begin{figure}
    \centering
    \begin{tikzpicture}
    \ctikzset{tripoles/mos style/arrows}[line width=thick, european]
´         % NMOS-Transistor-Symbol auf der linken Seite
        \begin{scope}[xshift=-3cm] % Verschiebung des Symbols nach links
            \draw (0,2) node[nmos](nmos){};
            \node[above] at (nmos.gate) {$G$}
                (nmos.gate) -- ++(-0.1,0) coordinate(gate)
                node[circle, draw,inner sep=1.5pt]{};
            \node[below, yshift=-5pt] at (nmos.source) {$S$}
                (nmos.source) -- ++(0,-0.1) coordinate(source)
                node[circle, draw,inner sep=1.5pt]{};
            \node[left] at (nmos.drain) {$D$}
                (nmos.drain) -- ++(0,0.1) coordinate(drain)
                node[circle, draw,inner sep=1.5pt]{};

            \draw[->, shorten >=5pt, shorten <=5pt, >=stealth] (drain) to[out=-45,in=45] (source)
                node[xshift=22pt, yshift=25pt]{\footnotesize{$V_{DS}$}};
            \draw[<-, shorten <=5pt, >=stealth] (drain) -- ++ (0,.75) 
                node[right, midway, yshift=5pt]{\footnotesize{$I_D$}};
            \draw[->, shorten >=5pt, shorten <=5pt, >=stealth] (gate) to[out=-90,in=180] (source)
                node[xshift=-25pt]{\footnotesize{$V_{GS}$}};
            \fill[white] (2,1) rectangle (2.975,0);
        \end{scope}

        \begin{axis}[
            xlabel={$V_{GS}$},
            ylabel={$I_D$},
            xmin=0, xmax=4,
            ymin=0, ymax=300,
            xtick=\empty, % Entfernt die Ticks und Beschriftungen der X-Achse
            ytick=\empty,
            width=7cm,
            height=6cm,
            %title={MOSFET Transferkennlinie},
            %title style={yshift=10pt}, % Abstand zwischen Titel und Grafik erhöhen
            %legend pos=north west,
            axis lines=middle, % Achsen in der Mitte mit Pfeilen
            xlabel style={at={(ticklabel* cs:1)}, anchor=north west, xshift=-10pt}, % Beschriftung neben    dem Pfeil
            ylabel style={at={(ticklabel* cs:1)}, anchor=south east, yshift=-10pt}, % Beschriftung neben    dem Pfeil
            minor tick num=0, % Schrittweite entfernen
            minor y tick num=0,
            extra x ticks={1.5}, % Markierung bei V_GS = 1.5
            extra x tick labels={$V_{th}$}, % Beschriftung der Markierung
            extra x tick style={
                grid=none, % Keine Gitterlinie bei der Markierung
                tick pos=left % Tick auf der linken Seite der Achse
            }
        ]
        
        % Beispiel: Quadratische Kennlinie im Sättigungsbereich
        \addplot[
            domain=1.5:4, % V_GS von 1V bis 7V
            samples=100,
            color=blue,
            thick,
        ] {150 * (x - 1.5)^4}; 
        \addplot[
            domain=0:1.5,
            samples=100,
            color=blue,
            thick,
        ] {0};
    
    
        \addplot[
            domain=1.5:4, % V_GS von 1V bis 7V
            samples=100,
            color=red,
            thick,
            style=dashed
        ] {100 * (x - 1.5)^4};
    
        \addplot[
            domain=1.5:4, % V_GS von 1V bis 7V
            samples=100,
            color=red,
            thick,
            style=dashed
        ] {250 * (x - 1.5)^4};
        
        \end{axis}
    \end{tikzpicture}
    \caption[Symbol NMOS-Transistor mit Transferkennlinie]{\\Quelle: Eigene Darstellung Symbol NMOS-Transistor mit Transferkennlinie in Anlehnung an \cite[S. 268]{MicroelectronicCircuits}}
    \label{fig:MOSFET_Transferkennlinie}
\end{figure}